
The calibration-alignment divergence gap grows at $\beta \approx$ 0.143--0.160 $nat^{-1}$ of KL optimization budget --- sufficient to traverse more than half the normalized scale by KL = 8 nats, and consistent enough to replicate within 12\% across two independent model families.

Four contributions are reported: (1) the calibration-alignment divergence construct, named and quantified; (2) slopes with confidence --- $\beta$ = 0.1433 $nat^{-1}$ ($R^2$ = 0.958, p < $10^{-6})$ with both parametric CI [0.119, 0.168] and bootstrap CI [0.117, 0.177] strictly positive; (3) cross-dataset replication with slope ratio 1.116; (4) the normalized divergence gap metric as a cross-study comparison instrument.

Three predictions were tested. P1 --- significantly positive slope in Coste et al. data --- was SUPPORTED with HIGH confidence. P2 --- independent replication in Gao et al. data --- was SUPPORTED with MEDIUM confidence (bootstrap CI marginally overlaps zero; parametric CI positive; parametric p = 0.0025). P3 --- $AI\rightarrow Human$ coverage ratio R > 0.90 across published alignment benchmarks --- was INCONCLUSIVE and is deferred to future work.

Future directions grounded in these results: (a) raw data access from original authors for exact $\beta$ with digitization uncertainty propagation; (b) piecewise linear regression on Gao et al. data to characterize the pre-onset regime; (c) behavioral study linking evaluation calibration divergence to user over-reliance; (d) coverage ratio R computation across the bidirectional alignment survey corpus; (e) cross-scale replication at multiple RM sizes to test whether $\beta \approx$ 0.14--0.16 $nat^{-1}$ is scale-invariant.

A field that measures alignment by proxy scores alone, without tracking held-out human preference, systematically overestimates how aligned its models are as optimization proceeds.

